\section{第4章：刚体的转动（4.5练习题）}
\sourceproblem{选择1}{56}{62}
\problemtext 定轴匀变速转动时距轴 $r$ 的点，比较切向和法向加速度的变化。
\workings $a_\tau=r|\alpha|$ 恒定，$a_n=r\omega^2$ 随 $\omega=\omega_0+\alpha t$ 变化，选C。
\crosscheck 取 $r>0,\alpha\ne0$；轴上点两分量均为零。
\sourceproblem{选择2}{56}{62}
\problemtext 两力作用于定轴刚体，判断合力与轴向力矩的逻辑关系。
\workings 平行轴的力对该轴力矩必为零，所以原A所说的“可能为零”也成立；垂直轴的两力也可使轴向力矩抵消，B成立。C、D分别可用力偶、通过轴线的合力反驳。
\crosscheck 严格按可能性的逻辑，A和B均真；若原意是区分必然与可能，应将A改成“可能不为零”等明确措辞。
\sourceproblem{选择3}{56}{62}
\problemtext 判断刚体内力矩、作用反作用力矩及相等质量刚体运动的命题（选项续PDF63）。
\workings 在教材经典中心内力假设下，$\bm\tau_{12}+\bm\tau_{21}=(\bm r_1-\bm r_2)\times\bm F_{12}=0$，选B。内力矩不改变整体轴向角动量；等质量不保证等惯量。
\crosscheck 若只假设作用反作用力等大反向、不假设沿两作用点连线，则力矩和可能非零，B须注明中心力条件。
\sourceproblem{选择4}{57}{63}
\problemtext 光滑定轴刚体受多个合力为零的力，判断转速。
\workings 轴向力矩可为零或非零，$\alpha=\tau_z/J$，所以转速可能不变也可能改变，选D。
\crosscheck 力偶的合力为零但力矩非零。
\sourceproblem{选择5}{57}{63}
\problemtext 匀质杆由水平静止绕端轴下摆，比较角速度及角加速度大小。
\workings 设下摆角 $\phi$，$\alpha=3g\cos\phi/(2l)$、$\omega^2=3g\sin\phi/l$，故角速率增大而角加速度大小减小，选C。
\crosscheck 最低点重力矩为零但角速率最大。
\sourceproblem{选择6}{57}{63}
\problemtext 飞轮绕绳挂质量 $m$ 重物时角加速度为 $\alpha$，改用恒拉力 $2mg$，比较新角加速度。
\workings 半径 $R$、惯量 $J$ 时 $\alpha=mgR/(J+mR^2)$，新值 $\alpha'=2mgR/J>2\alpha$，选C。
\crosscheck 原绳张力小于 $mg$，不是直接施加重力大小的恒力。
\sourceproblem{选择7}{57}{63}
\problemtext 相同定滑轮A受恒力 $Mg$，B挂质量 $M$ 重物，比较角加速度。
\workings $\alpha_A=MgR/J$，$\alpha_B=MgR/(J+MR^2)$，故 $\alpha_A>\alpha_B$，选B。
\crosscheck 理想无滑动及无轴摩擦。
\sourceproblem{选择8}{57}{63}
\problemtext 转动惯量取决于哪些因素。
\workings $J=\int r_\perp^2\,\mathrm dm$，取决于质量、分布、转轴位置和方向，选C。
\crosscheck 与转速无关。
\sourceproblem{选择9}{57}{63}
\problemtext 同半径同质量细环，一个质量均匀一个不均匀，比较中心垂直轴惯量。
\workings 各质量元轴距均为 $R$，故两者均为 $mR^2$，选C。
\crosscheck 转轴通过几何环心，不一定通过非均匀环质心。
\sourceproblem{选择10}{57}{63}
\problemtext 长 $l$、质量 $m$ 的匀质杆两端各有质量 $m$ 的小球，求中心垂直轴惯量。
\workings $J=ml^2/12+2m(l/2)^2=7ml^2/12$，选A。
\crosscheck 小球视作质点，系统质心仍位于中点。
\sourceproblem{选择11}{57}{63}
\problemtext 等质量等厚度两圆盘，$\rho_A>\rho_B$，比较半径与惯量（选项续PDF64）。
\workings $m=\rho\pi R^2h$，故 $R_A<R_B$；$J=mR^2/2$，故 $J_A<J_B$，选C。
\crosscheck 固定质量和厚度时惯量与密度成反比。
\sourceproblem{选择12}{58}{64}
\problemtext 判断角动量守恒的充分必要条件。
\workings 在惯性系取固定参考点，$\mathrm d\bm L/\mathrm dt=\bm\tau_{\rm ext}$，故合外力矩为零是充要条件，选B。
\crosscheck 若仅讨论固定轴分量，条件应是合外力矩的轴向分量为零，而非全部矢量必须为零。
\sourceproblem{选择13}{58}{64}
\problemtext 质点受有心力，判断何种力学量守恒。
\workings 对力心的力矩为零，角动量守恒，选C。
\crosscheck 有心力未必是时间无关的保守力，动能及势能也通常变化。
\sourceproblem{选择14}{58}{64}
\problemtext 自由转台上的人收回哑铃，判断机械能及角动量。
\workings 若把转台、人、哑铃整体计入，轴向外力矩为零，角动量守恒；惯量减小使 $K=L_z^2/(2J)$ 增加，机械能不守恒，选C为教材预期。
\crosscheck 原题只称人与哑铃系统，若转台具有惯量且可交换角动量，则该子系统角动量不一定守恒；严格表述应包含转台或说明其惯量可忽略。
\sourceproblem{选择15}{58}{64}
\problemtext 光滑铰链悬杆绕竖轴转动，小环沿杆下滑，判断杆环系统轴向角动量及含地球系统机械能。
\workings 若初始转动后系统自由演化，固定铰点不做功、重力对竖轴力矩为零，故机械能与轴向角动量均守恒，选D。
\crosscheck 若“匀角速”指外部驱动始终维持角速度恒定，则驱动可施加外力矩并做功，上述守恒不再自动成立；答案D采用自由演化解释。
\sourceproblem{选择16}{58}{64}
\problemtext 图示转盘受两个等大反向非共线力，判断转速变化。
\workings 图中原角速度为顺时针，力偶上方作用线向右、下方作用线向左，产生顺时针力矩，故角速率增大，选A。
\crosscheck 必须按两条力的作用线而非合力判定；图示方向保留复核标记。
\sourceproblem{选择17}{58}{64}
\problemtext 两等质量入射质点速度等大反向共线，分别嵌入转盘，比较碰后角速率。
\workings 两入射质点的轴向角动量之和为零，而碰后惯量增加，故由角动量守恒可知角速率减小，选C。
\crosscheck 原转速非零，嵌入位置距轴非零；不使用机械能守恒。
\sourceproblem{选择18}{58}{64}
\problemtext 自由定轴转动物体发生热膨胀或冷收缩，判断角速率（选项续PDF65）。
\workings 在通常均匀各向同性变形模型下，膨胀使轴距增大、惯量增大，角速率减小；收缩相反，选D。
\crosscheck 若变形仅沿轴向，惯量可不变，因此须注明径向尺寸变化的条件。
\sourceproblem{选择19}{59}{65}
\problemtext 惯量 $J$ 的水平转台初角速度 $\omega_0$，质量 $m$ 的人由中心沿径向到半径 $R$ 的边缘，求末角速度。
\workings 以人为质点，轴向角动量守恒，$\omega=J\omega_0/(J+mR^2)$，选A。
\crosscheck 假设人在转台系无额外切向相对速度；初末角动量均需包括人。
\sourceproblem{选择20}{59}{65}
\problemtext 匀质细杆长 $l$，由自然下垂位置获得初角速率，求持续绕端轴转动的条件。
\workings 质心到倒立位置升高 $l$，能量阈值由 $ml^2\omega^2/6=mgl$ 得 $\omega_c=\sqrt{6g/l}$。严格连续越过倒立点须 $\omega>\omega_c$，原选项C为阈值表达。
\crosscheck 等号时倒立点速度趋零，理想模型中趋近该点所需时间发散，不属于持续整周转动。
\sourceproblem{选择21}{59}{65}
\problemtext 长 $2l$、质量 $m$ 的细杆中心固定，质量 $m$ 小球以速率 $v$ 垂直撞击杆端并粘合，求末角速度。
\workings $mvl=(ml^2/3+ml^2)\omega$，故 $\omega=3v/(4l)$，选C。
\crosscheck 轴向角动量守恒，碰撞机械能损失为正。
\sourceproblem{选择22}{59}{65}
\problemtext 轻杆长20厘米，两球初始距中心各5厘米，烧断约束线后滑至杆端，求角速度。
\workings 无轴向外力矩，$J_i\omega_0=J_f\omega$，半径增为2倍，惯量变4倍，故 $\omega=\omega_0/4$，选D。
\crosscheck 角动量守恒不表示角速度恒定；此为球刚到杆端前的瞬时值。
\sourceproblem{选择23}{59}{65}
\problemtext 长 $5/3$ 米的匀质杆从与向下竖直成60度处静止释放，取 $g=10$，求最大角速率。
\workings $\tfrac12(ml^2/3)\omega^2=mg(l/2)(1-\cos60^\circ)$，故 $\omega=\sqrt{3g/(2l)}=3\,\mathrm{rad/s}$，选A。
\crosscheck 最低点势能最小；杆长为原页分数 $5/3$ 米。
\sourceproblem{选择24}{59}{65}
\problemtext 长 $L$、质量 $M$ 的杆绕端轴静止，质量 $m$ 的入射质点垂直穿过自由端，速度由 $v$ 降为 $v/2$，求杆角速率（续PDF66）。
\workings 短时轴向角动量守恒，$mvL-m(v/2)L=(ML^2/3)\omega$，得 $\omega=3mv/(2ML)$，选C。
\crosscheck 假设穿出方向不变；不使用机械能守恒。
\sourceproblem{选择25}{60}{66}
\problemtext 自转者收臂使转动惯量降为原来的三分之一，求角速度及动能。
\workings $J_0\omega_0=(J_0/3)\omega$，故 $\omega=3\omega_0$；$K=L_z^2/(2J)$，故 $K=3K_0$，选D。
\crosscheck 机械能增加来自内部做功。
\sourceproblem{选择26}{60}{66}
\problemtext 同长同质量单摆和端轴匀质杆由相同偏角静止释放，比较最低点角速率。
\workings 单摆 $\omega_1^2=2g(1-\cos\theta)/l$；杆 $\omega_2^2=3g(1-\cos\theta)/l$，故非零偏角时 $\omega_1<\omega_2$，选C。
\crosscheck 两者质心高度差和转动惯量不同；$\theta=0$ 时均为零。
\sourceproblem{选择27}{60}{66}
\problemtext 两同质量同平动速率球，一个无自转、另一个有自转，比较将其完全停止所做的功。
\workings 若指使两球平动和转动均停止的合外力总功，则 $W_1=-mv^2/2$，$W_2=-mv^2/2-J\omega^2/2$，故 $W_1>W_2$，选A。若比较克服运动的功的绝对值，则反向不等。
\crosscheck 若只停止质心而不停止自转，或只问手的功而另有耗散路径，须另给模型。
\sourceproblem{选择28}{60}{66}
\problemtext 惯量 $J$ 的圆盘角速率由 $\omega_0$ 降为 $\omega_0/2$，求阻力矩做功。
\workings 若阻力矩为唯一做功的力矩，$W=\Delta K=J(\omega_0^2/4-\omega_0^2)/2=-3J\omega_0^2/8$，选B。
\crosscheck 阻力矩做负功；有驱动力矩时不能仅由动能差确定阻力功。
\sourceproblem{选择29}{60}{66}
\problemtext 椭圆轨道卫星在远地点A和近地点B，比较角动量及动能。
\workings 中心引力对地心力矩为零，故 $L_A=L_B$；近远地点速度垂直径矢，$R_Av_A=R_Bv_B$，$R_A>R_B$ 时 $K_B>K_A$，选A。
\crosscheck 圆轨道退化情形两动能相等；题设最大最小距离不同。

\sourceproblem{填空1}{61}{67}
\problemtext 恒角加速度圆盘由 $20\pi$ 增至 $30\pi$ rad/s，期间转60圈，求角加速度及历时。
\workings $\Delta\theta=120\pi$，$\alpha=(\omega_2^2-\omega_1^2)/(2\Delta\theta)=25\pi/12\,\mathrm{rad/s^2}$，$\Delta t=(\omega_2-\omega_1)/\alpha=4.8\,\mathrm s$。
\crosscheck 平均角速率 $25\pi$ 乘4.8秒为 $120\pi$。
\sourceproblem{填空2}{61}{67}
\problemtext 飞轮初角速率5 rad/s，20秒后降为初值0.8倍，求恒角加速度及100秒转角。
\workings $\alpha=(4-5)/20=-0.05\,\mathrm{rad/s^2}$，$\theta=5(100)-0.025(100)^2=250\,\mathrm{rad}$。
\crosscheck 100秒恰好停止，区间内没有反向转动。
\sourceproblem{填空3}{61}{67}
\problemtext 皮带传动主动轮半径0.1米、1450转每分，从动轮半径0.29米，求从动轮缘速率和转速。
\workings 不打滑时 $v=2\pi(0.1)(1450/60)\approx15.18\,\mathrm{m/s}$；$n_2=1450(0.1/0.29)=500\,\mathrm{r/min}$。
\crosscheck 两轮缘线速率相同，角速率不同。
\sourceproblem{填空4}{61}{67}
\problemtext $\theta=10+8t-5t^2$（SI），求0.2秒角速度、角加速度及距轴0.5米点的速度和切法向加速度。
\workings $\omega=8-10t=6$，$\alpha=-10$（SI），$v=3\,\mathrm{m/s}$，$a_\tau=-5\,\mathrm{m/s^2}$（以正转向为正），$a_n=18\,\mathrm{m/s^2}$。
\crosscheck 切向加速度大小为5，法向指向轴。
\sourceproblem{填空5}{61}{67}
\problemtext $\omega=2+t+3t^2$，求0.2秒角速度及角加速度。
\workings $\omega=2.32\,\mathrm{rad/s}$，$\alpha=1+6t=2.2\,\mathrm{rad/s^2}$。
\crosscheck 角加速度是角速度的一阶导数。
\sourceproblem{填空6}{61}{67}
\problemtext 半径0.5米滑轮牵引升降机以0.4米每平方秒加速上升，求角加速度大小。
\workings 无滑动时 $|\alpha|=a/r=0.8\,\mathrm{rad/s^2}$。
\crosscheck 方向由绕绳方向确定。
\sourceproblem{填空7}{61}{67}
\problemtext 飞轮600转每分、惯量2.5，恒力矩制动1秒停止，求力矩大小。
\workings $\omega_0=20\pi$，$|M|=J\omega_0/t=50\pi\,\mathrm{N\,m}$。
\crosscheck 制动力矩方向与初始角速度相反。
\sourceproblem{填空8}{61}{67}
\problemtext 惯量2.0的轮受恒力矩负7.0牛米，8秒后角速度为初值的相反数，求初角速度。
\workings $\alpha=-3.5$，$-\omega_0=\omega_0-3.5(8)$，故 $\omega_0=14\,\mathrm{rad/s}$。
\crosscheck 4秒停止后反向加速，8秒末角速度为负14。
\sourceproblem{填空9}{61}{67}
\problemtext 惯量50的系统从静止0.5秒匀加速至 $4\pi$ rad/s，求0.3秒角速率及驱动力矩。
\workings $\alpha=8\pi$，$\omega(0.3)=2.4\pi\,\mathrm{rad/s}$，$M=400\pi\,\mathrm{N\,m}$。
\crosscheck 驱动力矩等于合力矩须忽略其他阻力矩。
\sourceproblem{填空10}{61}{67}
\problemtext 球惯量 $5\times10^{-2}$，$\theta=2+4t-t^2$，求首次停止时间和外力矩大小。
\workings $\omega=4-2t$，故 $t=2\,\mathrm s$；$\alpha=-2$，$|M|=0.10\,\mathrm{N\,m}$。
\crosscheck 角速度零不表示角加速度零。
\sourceproblem{填空11}{61}{67}
\problemtext 质量 $m$ 的斜抛质点从O出发，初速 $v_0$、仰角 $\alpha$，求对O角动量、导数及力矩。
\workings 取纸外为正，$x=v_0t\cos\alpha,\ y=v_0t\sin\alpha-gt^2/2$。故
\[
L_z=m(xv_y-yv_x)=-mgv_0t^2\cos\alpha/2,\qquad
\dot L_z=M_z=-mgv_0t\cos\alpha.
\]
\crosscheck 直接计算重力矩 $-mgx$ 与导数一致。
\sourceproblem{填空12}{61}{67}
\problemtext 同质量 $m$ 同半径 $R$ 的细环与均匀圆盘，求中心垂直轴惯量（续PDF68）。
\workings 细环 $J=mR^2$，圆盘 $J=mR^2/2$。
\crosscheck 圆盘各质量元轴距小于等于 $R$，因此惯量较小。
\sourceproblem{填空13}{62}{68}
\problemtext 两等质量等身高者从静止沿轻绳两端攀爬，甲相对绳速度为乙两倍，比较到顶时间。
\workings 理想轻滑轮两端张力相等，两人质量相同，故地面系加速度、速度及高度始终相同（相同初高条件下），同时到达。
\crosscheck 相对绳速度不同由绳的运动补偿；若初高度不同，不能推出同时到达。
\sourceproblem{填空14}{62}{68}
\problemtext 500千克飞轮质量集中在半径0.6米轮缘，角速率220 rad/s，30秒匀减速停止，求惯量和制动力矩。
\workings $J=500(0.6)^2=180\,\mathrm{kg\,m^2}$，$|M|=180(220/30)=1320\,\mathrm{N\,m}$。
\crosscheck 质量集中于轮缘，不使用均匀圆盘的二分之一系数。
\sourceproblem{填空15}{62}{68}
\problemtext 飞轮20秒由40 rad/s匀减速至零，惯量50，求转角及力矩大小。
\workings $\theta=(40/2)20=400\,\mathrm{rad}$，$\alpha=-2$，$|M|=100\,\mathrm{N\,m}$。
\crosscheck 角速度时间图面积为400。
\sourceproblem{填空16}{62}{68}
\problemtext 长 $L$ 轻杆两端固定 $m,2m$，绕中点水平轴转动，从图示60度释放，求惯量、水平位置力矩及角加速度大小。
\workings $J=3mL^2/4$。水平位置重力矩大小 $(2m-m)gL/2=mgL/2$，故 $|\alpha|=2g/(3L)$。
\crosscheck 重球一侧下摆；两球重力矩方向相反。
\sourceproblem{填空17}{62}{68}
\problemtext 匀质杆长 $L$，水平静止释放，求初角速度和初角加速度。
\workings $\omega_0=0$，$J=mL^2/3$，$|\alpha_0|=mgL/(2J)=3g/(2L)$。
\crosscheck 初始角速度零而重力矩非零。
\sourceproblem{填空18}{62}{68}
\problemtext 匀质细杆长 $l$、质量 $m$，绕端轴由水平静止下摆角 $\theta$，求惯量和角速率。
\workings $J=ml^2/3$，$\frac12J\omega^2=mg(l/2)\sin\theta$，故 $\omega=\sqrt{3g\sin\theta/l}$。
\crosscheck $\theta=0$ 为初态，$\theta=\pi/2$ 为最低点。
\sourceproblem{填空19}{62}{68}
\problemtext 匀质杆水平静止释放，求初角加速度及竖直位置角速率。
\workings $\alpha_0=3g/(2l)$（下摆为正），$\omega_{\rm bottom}=\sqrt{3g/l}$。
\crosscheck 分别使用转动定律和机械能守恒。
\sourceproblem{填空20}{62}{68}
\problemtext 匀质杆长 $L$ 由竖直直立轻碰倒下，底端固定不滑动，求将达水平地面时角速率及角加速度大小。
\workings 取微扰极限，质心下降 $L/2$，$\omega=\sqrt{3g/L}$；水平位置 $|\alpha|=3g/(2L)$。
\crosscheck 精确直立且静止为不稳定平衡；结果为碰地前瞬时值，不涉及碰撞。
\sourceproblem{填空21}{62}{68}
\problemtext 匀质杆AB长 $l$，由水平静止下摆至角 $\theta$，求B端速率及重力矩功。
\workings $W=mg(l/2)\sin\theta$，$\omega^2=3g\sin\theta/l$，故 $v_B=\sqrt{3gl\sin\theta}$。
\crosscheck 不能将端点下降高度当作质心下降高度。
\sourceproblem{填空22}{62}{68}
\problemtext 恒轴向力矩 $M$ 作用于角速度 $\omega$ 的刚体，求瞬时功率及之后转两圈的功。
\workings 采用有向量，$P=M\omega$，若沿正方向转两圈，$W=4\pi M$。
\crosscheck 若转向与力矩反向，功为负；应采用有向转角。
\sourceproblem{填空23}{62}{68}
\problemtext 惯量 $J$ 的刚体受恒合力矩 $M$，求角加速度及角速度为 $\omega$ 时的功率。
\workings $\alpha=M/J$，$P=M\omega$。
\crosscheck 量纲分别为角度每平方秒和瓦。
\sourceproblem{填空24}{62}{68}
\problemtext 直径 $D$、质量 $m$ 的均匀飞轮从静止经时间 $t$ 被恒绳力加速至角速率 $\omega$，求角加速度及拉力功（续PDF69）。
\workings $\alpha=\omega/t$，$J=mD^2/8$，若忽略其他力矩做功，$W=J\omega^2/2=mD^2\omega^2/16$。
\crosscheck 也可由恒力矩 $J\alpha$ 乘转角 $\omega t/2$ 得到。
\sourceproblem{填空25}{63}{69}
\problemtext 地球两极冰融化后，判断地球惯量与自转速率的变化。
\workings 教材模型中质量由近轴区域向较大轴距分布，惯量增大；近似无外轴向力矩时 $\omega=L_z/J$ 减小。
\crosscheck 实际变化取决于水质量重新分布、地球形变及其他角动量交换，非仅由融化本身决定。
\sourceproblem{填空26}{63}{69}
\problemtext 惯量 $J$ 的圆盘受阻力矩 $-k\omega$，由 $\omega_0$ 减至一半，求时间和阻力矩功。
\workings $J\dot\omega=-k\omega$，$\omega=\omega_0e^{-kt/J}$，故 $t=J\ln2/k$，$W=-3J\omega_0^2/8$。
\crosscheck $k>0$；阻力矩为唯一做功项。
\sourceproblem{填空27}{63}{69}
\problemtext 惯量 $J_0$ 的飞轮以 $\omega_0$ 转动，与静止、惯量为其两倍的飞轮接合，求末角速度与动能变化。
\workings $\omega=\omega_0/3$，$\Delta K=(3J_0)(\omega_0/3)^2/2-J_0\omega_0^2/2=-J_0\omega_0^2/3$。
\crosscheck 轴向角动量守恒，机械能不守恒。
\sourceproblem{填空28}{63}{69}
\problemtext 光滑桌面小球质量 $m$，初始圆周半径 $r_1$、角速度 $\omega_1$，缓慢拉绳缩至 $r_2$，求动能增量。
\workings 角动量 $L=mr_1^2\omega_1$ 守恒，初末均为圆周运动时
\[
\Delta K=\frac{m r_1^2\omega_1^2}{2}\left(\frac{r_1^2}{r_2^2}-1\right).
\]
\crosscheck $r_2<r_1$ 时动能增加，由拉力做正功；若末态存在径向速度，须计入其动能。
\sourceproblem{填空29}{63}{69}
\problemtext 滑冰者惯量由 $J_0$ 减至 $2J_0/3$，初角速度 $\omega_0$，求末角速度及动能变化。
\workings $\omega=3\omega_0/2$，$K_f=(3/2)K_i$，故 $\Delta K=J_0\omega_0^2/4$。
\crosscheck 角动量守恒；增加的能量来自人体内部做功。
\sourceproblem{填空30}{63}{69}
\problemtext 长 $L$、质量 $M$ 的均匀杆由水平静止下摆，在竖直位置与静止质量 $m$ 物体完全非弹性碰撞，求碰后物体速率及对轴角动量。
\workings 碰前 $\omega_i=\sqrt{3g/L}$。若物体与杆端粘合，短时轴向角动量守恒给出
\[
\omega_f=\frac{M}{M+3m}\omega_i,\quad
v_m=\frac{M\sqrt{3gL}}{M+3m},\quad
|L_m|=mLv_m.
\]
\crosscheck 假设地面水平冲量可忽略、碰撞点在杆端；若物体未粘合或有其他外冲量矩，需重新建模。

\sourceproblem{计算1}{63}{69}
\problemtext 飞轮从静止6.0秒内匀加速至36 rad/s，求角加速度和转过圈数。
\workings $\alpha=36/6=6\,\mathrm{rad/s^2}$，$\theta=36(6)/2=108$ rad，$N=54/\pi\approx17.19$ 圈。
\crosscheck 角速度时间图面积为108 rad，不能将弧度直接当圈数。
\sourceproblem{计算2}{63}{69}
\problemtext 定轴运动 $\theta=at+bt^2-ct^4$，求角加速度及轴距 $r$ 点的切、法向加速度。
\workings $\omega=a+2bt-4ct^3,\alpha=2b-12ct^2$，$a_\tau=r(2b-12ct^2)$ 为沿正转向有向分量，$a_n=r(a+2bt-4ct^3)^2$。
\crosscheck 法向分量非负且指向轴；切向大小取绝对值。
\sourceproblem{计算3}{64}{70}
\problemtext 齿轮角加速度 $\alpha=4at^3-3bt^2$，初角速度 $\omega_0$，求角速度及转角。
\workings 积分得 $\omega=\omega_0+at^4-bt^3$，$\Delta\theta=\omega_0t+at^5/5-bt^4/4$。
\crosscheck 两次微分恢复给定角加速度；若发生反转，实际转过的总角路程须积分 $|\omega|$。
\sourceproblem{计算4}{64}{70}
\problemtext 均匀圆板半径 $R$、质量 $m$，在水平粗糙桌面绕中心轴转动，动摩擦系数 $\mu$，接触压力均匀，求摩擦力矩。
\workings 面压力 $p=mg/(\pi R^2)$，半径 $r$ 的环带摩擦力大小 $\mathrm dF=\mu p\,2\pi r\,\mathrm dr$，故
\[
|M_f|=\int_0^R r\,\mathrm dF=\frac23\mu mgR.
\]
方向与角速度相反。
\crosscheck 有效摩擦力臂为 $2R/3$，不是圆板质心的轴距零，也不是边缘半径。
\sourceproblem{计算5}{64}{70}
\problemtext 长 $l$、质量 $M$ 的均匀细杆绕端轴，另一端固定质量 $m$、半径 $r$ 的均匀实心小球，求系统惯量。
\workings 若杆端连接球面且球心在杆延长线上，轴到球心距 $l+r$，由平行轴定理
\[
J=Ml^2/3+2mr^2/5+m(l+r)^2.
\]
若 $l$ 指轴到球心的距离，则最后一项改为 $ml^2$，杆自身长度按实际取值。
\crosscheck 原图连接点与球心距离须区分；球质量分布若非均匀实心球，应另给质心惯量。
\sourceproblem{计算6}{64}{70}
\problemtext 圆柱半径 $R$，有四个贯通圆孔，孔半径 $R/3$，孔心距主轴 $R/2$，求中心轴惯量。
\workings 设未挖孔圆柱质量 $M_0$，每孔等效质量 $M_0/9$。若四孔互不重叠，则减去各孔惯量
\[
J=\frac12 M_0R^2-4\frac{M_0}{9}\left(\frac{R^2}{18}+\frac{R^2}{4}\right)=\frac{59}{162}M_0R^2.
\]
若题给 $m$ 是挖孔后质量，$m=5M_0/9$，则 $J=59mR^2/90$；若指完整圆柱质量，则取 $M_0=m$。
\crosscheck 相邻孔心距离 $R/\sqrt2>2R/3$，各孔不交叠；最外缘 $5R/6<R$，减法成立。原题质量含义需明示。
\sourceproblem{计算7}{65}{71}
\problemtext 实心圆柱质量16 kg、半径15 cm，轻绳悬8 kg物体，取 $g=10$，求张力及静止释放1秒下落距离。
\workings $J=Mr^2/2$，$mg-T=ma,\ T=Ja/r^2$，得 $a=mg/(m+M/2)=5\,\mathrm{m/s^2}$，$T=40$ N，$s=at^2/2=2.5$ m。
\crosscheck 张力小于悬物重力；半径在无滑动实心柱模型中约去。
\sourceproblem{计算8}{65}{71}
\problemtext 两质量 $m_1,m_2$ 分别悬在半径 $r,R$ 的同轴无质量鼓轮上，绳绕向相反，求角加速度。
\workings 取 $m_2$ 下降方向为正，$a_2=R\alpha,a_1=r\alpha$ 向上。消去两绳张力得
\[
\alpha=\frac{(m_2R-m_1r)g}{m_1r^2+m_2R^2}.
\]
\crosscheck 驱动力矩差为零时角加速度为零；若分子为负实际方向反转。
\sourceproblem{计算9}{65}{71}
\problemtext 光滑斜面倾角 $\theta$，物体质量 $m$ 经轻绳带动质量 $M$、半径 $R$ 的固定轴滑轮转动，求下滑加速度。
\workings $mg\sin\theta-T=ma,\ TR=Ja/R$，故 $a=mg\sin\theta/(m+J/R^2)$。若滑轮为均匀实心圆盘，$J=MR^2/2$，得 $a=mg\sin\theta/(m+M/2)$。
\crosscheck 题面仅给质量和半径，不足以唯一确定转动惯量；圆盘假设须明确。
\sourceproblem{计算10}{66}{72}
\problemtext 斜面30度、物体6 kg、悬物18 kg，均匀圆盘滑轮2 kg，斜面摩擦系数0.1，取 $g=10$，求加速度与两侧张力。
\workings 设悬物下降、斜面物体上行，$f=0.1(6)(10)\cos30^\circ=3\sqrt3$ N。
\[
a=\frac{180-30-3\sqrt3}{18+6+1}=\frac{150-3\sqrt3}{25}\approx5.792\,\mathrm{m/s^2}.
\]
斜面侧 $T_1=6a+30+3\sqrt3\approx69.95$ N；悬物侧 $T_2=180-18a\approx75.74$ N。
\crosscheck $T_2-T_1=a$（SI数值式），符合滑轮转动方程；取题给摩擦系数为滑动摩擦系数，启动需克服静摩擦。
\sourceproblem{计算11}{66}{72}
\problemtext 均匀圆盘原质量 $M$、半径 $R$，绕水平中心轴以角速度 $\omega$ 转动，边缘质量 $m$ 的碎块脱离并竖直上抛，求上升高度及余盘角动量。
\workings 若脱离无额外冲量，碎块初速 $R\omega$，相对脱离点最高上升 $h=R^2\omega^2/(2g)$。原盘角动量 $MR^2\omega/2$，碎块脱离瞬间角动量 $mR^2\omega$，故余盘 $L=(M/2-m)R^2\omega$。
\crosscheck 假设碎块尺寸可忽略且轴向外冲量矩为零；最高点绝对高度还需脱离点高度。
\sourceproblem{计算12}{66}{72}
\problemtext 匀质杆长 $l$，从水平上方60度静止释放，绕端轴下摆，$g=10$，求初角加速度与过水平时角速率。
\workings 下摆为正，初力矩 $mg(l/2)\cos60^\circ$，故 $\alpha_0=3g/(4l)=7.5/l$。质心到水平下降 $(l/2)\sin60^\circ$，能量守恒得 $\omega=\sqrt{3g\sin60^\circ/l}=\sqrt{15\sqrt3/l}$。
\crosscheck 初角加速度并非水平位置的 $3g/(2l)$；角速度由势能差求得。
\sourceproblem{计算13}{66}{72}
\problemtext 匀质水平杆质量 $m$、长 $l$，两端竖直细线悬挂，烧断一根，求另一根瞬时张力。
\workings 设保留端为A。初速度为零且细线不可伸长，A瞬时竖直加速度为零；杆此刻角加速度可按绕A求出，$mg(l/2)=(ml^2/3)\alpha$，故 $\alpha=3g/(2l)$。质心向下加速度 $\alpha l/2=3g/4$，因此 $mg-T=3mg/4$，$T=mg/4$。
\crosscheck 也可用质心转动方程 $Tl/2=(ml^2/12)\alpha$ 联立；A并非全过程固定铰点，此处仅用瞬时约束。
